\section{The instrument and an explicit finite pulse choice}
\label{sec:chvisd-instrument}

Write $P_1=P_{ba}$ and $P_0=I_{16}-P_{ba}$. In the ideal frozen-control
model with a blank pointer, the joint source--field Kraus operators and
the measurement channel are
\begin{equation}
 \begin{aligned}
 K_y&=(P_yR)\otimes I_f,\\
 \mathcal J_y(\rho)&=K_y\rho K_y^*,\\
 \sum_yK_y^*K_y&=I,\\
 \mathcal J(\rho)&=\sum_{y=0}^1|y\rangle\langle y|
                         \otimes\mathcal J_y(\rho).
 \end{aligned}
\label{eq:chvisd-ideal-instrument}
\end{equation}
Thus $K_y^*K_y=E_y\otimes I_f$. The post-readout source can be disturbed
by the projector; the field is retained in each output rather than
replaced by a scalar marker. With the actual full drift and pulses,
$\widetilde K_y=\langle y|U_{\rm gate}U_{\rm rot}|0\rangle$
is in general a joint source--field operator. Unitarity still gives
$\sum_y\widetilde K_y^*\widetilde K_y=I$.
The complete ideal and actual postmeasurement channels, with outcomes
and quantum output retained, are within trace distance $r$ on the
reference input because the corresponding premeasurement pure states
are within vector distance $r$. Applying the same channel to a
preparation within trace distance $\delta_{\rm prep}$ adds at most
$\delta_{\rm prep}$ to this joint-output distance. The two instruments
therefore carry an all-outcome bound, not only a contrast estimate.

Use rectangular nominal pulses of equal finite duration
\begin{equation}
\begin{gathered}
\tau_{\rm rot}=\tau_{\rm gate}=\tau=\frac1{32768},\\
 C^{\rm up}=2\epsilon^{\rm up}+6c^{\rm up}
                   +4K_{\rm prep}^{\rm up}+4\epsilon^{\rm up}E_* .
\end{gathered}
\label{eq:chvisd-explicit-times}
\end{equation}
The exact native comparison evaluates
\begin{equation}
 r^2\leq2C^{\rm up}\tau\left(2+\frac{5\pi_*}{2}\right)^2
       \leq6.7949588\times10^{-7}<10^{-6}
\label{eq:chvisd-pulse-feasibility}
\end{equation}
for all six lanes. In particular $r\leq0.00082431540$ across the
six lanes, leaving the declared $10^{-3}$ control allowance sufficient
for the nominal mathematical pulses. Their strengths and absolute
switching-work ceiling are finite and explicitly charged:
\begin{equation}
\begin{aligned}
 |\lambda_{\rm rot}|&\leq\frac{\pi_*}{4\tau}
                                      <25735.928,\\
 |\lambda_{\rm gate}|&\leq\frac{\pi_*}{\tau}
                                      <102943.71,\\
 W_{\rm abs}^{\rm switch}&\leq\frac{5\pi_*}{2\tau}
                                      <257359.28.
\end{aligned}
\label{eq:chvisd-pulse-resources}
\end{equation}
These are original model units, without an assigned SI conversion or
an observed control calibration. They display the strength/work cost
of this short-pulse construction. An actuator/reference implementing
the spectral projector and exchange drive must supply its charge,
phase, timing and energy resources. The two pulses take the positive
latency $1/16384$ before subsequent pointer registration. The ideal
observable remains the observable at pulse entry.


